% ==================================================================
% id:       func-027
% topic:    מספרים ממשיים
% subtopic: אי־שוויונות וערך מוחלט
% solution: full
% ==================================================================
\textbf{יהיו $x,y\in\mathbb{R}$.}
\begin{enumerate}
\item
הוכיחו כי
$-|x|\le x\le|x|$.
\begin{solbox}
נפריד לשני מקרים.
אם $x\ge0$, אז $x=|x|$, ובנוסף $-|x|\le0\le x$. לכן $-|x|\le x=|x|$.
אם $x<0$, אז $x=-|x|$, ובנוסף $x<0\le|x|$. לכן $-|x|=x<|x|$.
בשני המקרים מתקיים
$-|x|\le x\le|x|$.
\end{solbox}
\item
הוכיחו את אי-שוויון המשולש:
$|x+y|\le|x|+|y|$.
\hint{היעזרו בסעיף א׳.}
\begin{solbox}
לפי סעיף א מתקיים
$-|x|\le x\le|x|$
וגם
$-|y|\le y\le|y|$.
נחבר את שני אי-השוויונות ונקבל
$$-\left(|x|+|y|\right)\le x+y\le|x|+|y| .$$
לכל $a\ge0$ מתקיים $|t|\le a$ אם ורק אם $-a\le t\le a$. נציב $t=x+y$
ו-$a=|x|+|y|$ ונקבל
$|x+y|\le|x|+|y|$.
\end{solbox}
\item
הסיקו כי
$\bigl||x|-|y|\bigr|\le|x-y|$.
\begin{solbox}
לפי סעיף ב,
$|x|=|(x-y)+y|\le|x-y|+|y|$,
ולכן
$|x|-|y|\le|x-y|$.
בהחלפת התפקידים של $x$ ו-$y$:
$|y|-|x|\le|y-x|=|x-y|$.
שני אי-השוויונות יחד אומרים כי
$-|x-y|\le|x|-|y|\le|x-y|$,
כלומר
$\bigl||x|-|y|\bigr|\le|x-y|$.
\end{solbox}
\end{enumerate}
